%%%%%%%%%%%%%%%%%%%%%%%%%%%%%%%%%%%%%%%%%%%%%%%%%%%%%%%%%%%%%%%%%%%%%%%%%%%%%%%%%%%%%%%%%%%%%%%%%%% 
% KAPITOLA 1
%%%%%%%%%%%%%%%%%%%%%%%%%%%%%%%%%%%%%%%%%%%%%%%%%%%%%%%%%%%%%%%%%%%%%%%%%%%%%%%%%%%%%%%%%%%%%%%%%%
\chapter{Úvod}
\label{ch:kap1}
Testovanie je jednou z najdôležitejších fáz vývoja softvéru, pretože dôsledky nekvalitného softvéru sú nákladné. Nekvalitný softvér spôsobil v roku 2022 v USA náklady odhadované na približne 2,41 bilióna amerických dolárov~\cite{Krasner2022}. Jedným z prístupov k testovaniu je testovanie riadené dátami (angl. Data-Driven testing, DDT), testovanie založené na princípe oddelenia testovacích dát od testovacej logiky.


Efektívne testovanie stojí na kvalite testovacích dát, hlavne preto, že pri nedostatočne domyslených dátach môžu chyby ostať neodhalené, čo vedie k časovým aj finančným stratám.


Vytváranie testovacích dát je časovo náročné a vyžaduje znalosť štruktúry testovaného systému, ale aj charakteru reálnych dát.


Na uľahčenie práce existujú nástroje na generovanie náhodných testovacích dát. Jedným z nich je Dataster~\cite{Vlach2020}, nástroj, ktorý generuje štruktúrované testovacie dáta pre relačné databázy na základe konfiguračného receptu. Ručné vytvorenie receptu vo formáte YAML vyžaduje znalosť dokumentácie programu Dataster\footnote{\url{https://dataster.testos.org/}}, ako aj znalosť schémy databázy, vzťahov medzi tabuľkami a charakteru uložených hodnôt. Náročnosť tohto procesu rastie so zložitosťou databázy. Pre používateľa s obmedzenými znalosťami predstavuje tvorba receptu pre konkrétnu databázu značné množstvo práce.


Potenciálnym riešením je automatizácia tejto úlohy s pomocou veľkých jazykových modelov (LLM), ktoré majú schopnosť interpretovať štruktúrované dáta a transformovať ich do iných formátov~\cite{Naveed2024}. Táto práca využíva lokálne LLM, keďže vstupom sú reálne dáta z databázy a lokálny beh zabezpečuje nielen súkromie dát, ale aj nízke náklady.


Cieľom tejto práce je navrhnúť a implementovať nástroj, ktorý sa pripojí k relačnej databáze, extrahuje jej schému spolu so vzorkami uložených dát a pomocou lokálnych LLM automaticky vygeneruje recept pre nástroj Dataster. Ďalším cieľom je zhodnotiť kvalitu vygenerovaných receptov naprieč rôznymi modelmi a konfiguráciami. Kvalita sa meria od správnosti vygenerovaného receptu (YAML) cez akceptáciu nástrojom Dataster až po vykonateľnosť výsledného SQL kódu. Vyhodnotenie prebieha na sade testovacích databáz rôznej zložitosti.



Kapitola~\ref{ch:kap2} približuje definíciu testovania a popis DDT. Kapitola~\ref{ch:kap3} definuje LLM, ich využitie pri generovaní dát, predstavuje jednotlivé modely a alternatívne prístupy k prispôsobeniu modelov úlohe. Kapitola~\ref{ch:kap4} približuje nástroj Dataster, jeho využitie, limitácie, popis konfiguračného receptu a formátu YAML. Kapitola~\ref{ch:kap5} predstavuje vysokoúrovňový návrh programu. Kapitola~\ref{ch:kap6} sa zaoberá implementáciou. Kapitola~\ref{ch:kap7} vysvetľuje proces experimentov, popisuje sadu testovacích databáz, úrovne validácie vygenerovaných receptov a porovnávané modely a konfigurácie. Kapitola~\ref{ch:kap8} prezentuje a porovnáva výsledky experimentov. Kapitola~\ref{ch:kap9} rozoberá limitácie programu a možné rozšírenia v budúcej práci. Na záver Kapitola~\ref{ch:kap10} zhŕňa vykonanú prácu, dosiahnuté výsledky a porovnáva ich s cieľmi práce.
